\section{Discussion and extensions}

The characterization in \cref{obj:thm:log-free-bivariate-capacity} gives an interaction-count benchmark for experimental sample size. Under the independent uniform sampling law of \cref{obj:ass:uniform-draw}, the relevant comparison is between \(B=\binom d2\), the number of coordinate pairs, and \(n\), the number of experimental units. For the centered prognostic class with the pooled Sobolev budget in \cref{obj:def:sobolev-class}, minimax scaled excess variance is comparable to \(\min\{1,(B/n)^s\}\). Thus the interaction count summarizes the dimension dependence of the worst-case precision requirement within this class.

\subsection{Causal calibration and relative excess}

To interpret the assignment guarantee in terms of causal precision, we first specify the potential outcomes, treatment, estimator, and targets used in the Gaussian calibration.

\begin{definitionv}{def:causal-completion}[Causal completion and estimator \(\widehat\theta\)]
The actual estimator \(\widehat\theta\) uses \(Y\), the pre-assignment potential-outcome array, and the treatment indicator
\[
A=(Z+1)/2,
\]
where \(Z\) is the vector of assignment signs. Its conditional estimand is \(\theta_n\), the average realized potential-outcome contrast, conditional on the realized full schedule. Its sampling-law estimand is \(\theta\), the sampling-law mean of that contrast. The setting and constructions are specified in \cref{obj:ass:order-two,obj:ass:pooled-budget,obj:ass:uniform-draw,obj:ass:fair,obj:ass:assignment-independence,obj:def:sobolev-class,obj:def:design-class,obj:def:criterion,obj:def:cosine-prior,obj:def:capacity-handle,obj:def:ordered-boundary-rounding,obj:lem:cosine-prior,obj:thm:full-design-lower,obj:lem:reflection-transfer,obj:lem:ht-transfer,obj:prop:published-class-inclusion,obj:thm:log-free-bivariate-capacity,obj:lem:isotropic-row-signing,obj:lem:ordered-prefix-rounding,obj:lem:published-prognostic-capacity,obj:thm:capacity-answer,obj:lem:exact-reflection-budget,obj:synth_1,obj:synth_2,obj:synth_3,obj:synth_4,obj:synth_5}.
\end{definitionv}

Concretely, for each \(m\in\mathcal H_{d,s}\), the completion in \cref{obj:synth_4} fixes
\[
\tau(x)=0.2+0.1\sqrt2\sin(2\pi x_1)
\]
and generates the pre-assignment schedule as
\[
Y_i(0)=m(X_i)-\frac{\tau(X_i)}2+\epsilon_i(0),
\qquad
Y_i(1)=m(X_i)+\frac{\tau(X_i)}2+\epsilon_i(1).
\]
The \(2n\) noise variables are independent standard Gaussians, independent of the covariates. After this schedule is generated, the admissible assignment law draws \(Z\) using the covariates and satisfies \cref{obj:ass:fair,obj:ass:assignment-independence}. With \(A_i=(Z_i+1)/2\), the observed outcome is \(Y_i(A_i)\), and the estimator and its two targets are
\[
\widehat\theta=\frac2n\sum_{i=1}^n Z_iY_i(A_i),
\qquad
\theta_n=\frac1n\sum_{i=1}^n\bigl(Y_i(1)-Y_i(0)\bigr),
\qquad
\theta=\mathbb E[\theta_n]=0.2.
\]
Fairness and outcome-independent assignment make the estimator conditionally unbiased for the realized schedule average \(\theta_n\), and its sampling-law expectation is \(\theta\).

The benchmark for this family is \(V^*=4.01\): treatment-effect heterogeneity contributes \(0.01\), and the two unit conditional outcome variances contribute \(4\). Here variance is taken over covariate sampling, potential-outcome noise, and assignment. The exact identity in \cref{obj:lem:ht-transfer} is
\[
n\operatorname{Var}(\widehat\theta)-V^*
=
\mathcal L_{n,d,s}(\pi,m).
\]
Consequently, the signing criterion measures actual Horvitz--Thompson excess variance above \(V^*/n\) for this completion. The Gaussian family calibrates the abstract imbalance loss while realizing every prognostic function in the pooled class.

A relative-excess tolerance \(\varepsilon>0\) requires
\[
n\operatorname{Var}(\widehat\theta)
\leq V^*(1+\varepsilon)
\]
uniformly over the Gaussian completions indexed by \(m\in\mathcal H_{d,s}\). For the assignment procedure in \cref{obj:def:capacity-handle}, the sufficient sample-size implication of \cref{obj:thm:log-free-bivariate-capacity} is
\[
n\geq B\max\left\{
1,\left(\frac{3072}{4.01\varepsilon}\right)^{1/s}
\right\}
\quad\Longrightarrow\quad
\sup_{m\in\mathcal H_{d,s}}
\left\{
\frac{n\operatorname{Var}(\widehat\theta)}{4.01}-1
\right\}
\leq\varepsilon.
\]
The same assignment procedure supplies this guarantee for every \(0<s\leq1\).

Conversely, any admissible design meeting that uniform relative-excess requirement must have \(R_s(n,d)\leq4.01\varepsilon\). For \(0<\varepsilon<c_s/4.01\), where \(c_s\) is the lower comparison constant in \cref{obj:thm:log-free-bivariate-capacity}, this entails
\[
n\geq B\left(\frac{c_s}{4.01\varepsilon}\right)^{1/s}.
\]
The tolerance restriction places the necessary implication below the saturated lower bound. At \(s=1\), \(c_1\approx3.36\times10^{-7}\), hence \(c_1/4.01\approx8.4\times10^{-8}\), while \(3072/c_1>9\times10^9\). The constants therefore provide conservative finite-sample certificates of the \(B\varepsilon^{-1/s}\) scaling; they should not be read as a sharp numerical calibration.

\subsection{Precision over bounded-density classes}

The broader comparison in \cref{obj:thm:log-free-bivariate-capacity} gives a necessary precision condition beyond uniform sampling. The class \(\mathcal A_{d,s}\) defined in \cref{obj:lem:published-prognostic-capacity} allows covariate densities between fixed positive finite bounds containing one. Its potential outcomes have finite second moments, and its conditional midpoint mean is a permitted intercept plus a centered member of \(\mathcal H_{d,s}\), with squared midpoint mean integrating to at most one under the covariate law.

Because this class contains the uniform Gaussian completions, its minimax scaled excess satisfies
\[
R_s(n,d)
\leq
\inf_{\pi\in\mathcal D_{n,d,s}}
\sup_{P\in\mathcal A_{d,s}}\mathcal V_n(\pi,P),
\]
where \(\mathcal V_n(\pi,P)\) is excess above the law-specific benchmark specified in \cref{obj:thm:log-free-bivariate-capacity}. Thus a uniform precision requirement over the broader class must accommodate the uniform-sampling lower bound. In particular, for fixed \(s\) and fixed density bounds, vanishing minimax scaled excess over these classes requires \(d_n=o(\sqrt n)\). This is the necessary interaction-count condition inherited from the uniform subclass.

\subsection{Limitations and future work}

Attainment under nonuniform covariate densities is a natural next question. The two-sided characterization and the specified assignment guarantee apply to independent uniform sampling, while the bounded-density comparison supplies a converse. Extending attainment would require assignment and approximation bounds that account for the covariate density.

A second direction concerns prognostic functions beyond the exact centered order-two decomposition and pooled budget in \cref{obj:ass:order-two,obj:ass:pooled-budget}. Higher-order interactions, approximation error in the decomposition, and unknown prognostic intercepts call for explicit control of their contributions to imbalance. Such extensions could clarify how sample-size requirements change when pairwise structure is an approximation to the outcome regression.

Finally, the finite construction in \cref{obj:def:ordered-boundary-rounding} uses exact-real arithmetic conventions. Finite-precision certification and bit complexity remain implementation questions. Quantifying how numerical error affects boundary decisions, fairness, and row imbalance would connect the mathematical assignment guarantee to a certified numerical procedure.
